\section{Welfare}\label{sec:welfare}

The coordinated benchmark in Section~\ref{sec:equilibrium} internalizes route-level system returns but is not identical to social welfare. This section introduces a quasilinear microfoundation that separates the private, coordinated, and social effort choices and shows how the associated wedges evolve with downstream reallocation.

\subsection{A microfounded benchmark}

Let the masses of consumers attached to routes $S$ and $R$ be $1-h(\eta)$ and $h(\eta)$, respectively, where
\[
    0<h(\eta)<1, \qquad h'(\eta)>0.
\]
A representative consumer obtains quasilinear utility
\begin{equation}
    u(q,x)=(a+x)q-\frac{q^2}{2}-pq,
    \label{eq:utility}
\end{equation}
so per-capita demand is
\begin{equation}
    d(x)=a+x-p.
    \label{eq:demand}
\end{equation}
We assume $a>p>c_S>c_R\geq0$, where $c_S$ and $c_R$ are real distribution costs. Product improvement and effort cost are
\[
    x=\beta e, \qquad C(e)=\frac{k e^2}{2},
\]
with $\beta>0$ and $k>\beta^2$. Aggregate route quantities are
\[
    q_S=(1-h)d(x), \qquad q_R=h d(x).
\]

Let the information producer's private margin be $m$, with $0<m<p-c_S$. Writing the remaining upstream component of the route-$S$ system margin as $b>0$, we have $b+m=p-c_S$. The coordinated system margins are therefore $p-c_S$ on route $S$ and $p-c_R$ on route $R$.

\subsection{Private, coordinated, and social effort}

The private effort choice solves
\[
    \max_e\; m(1-h)(a+\beta e-p)-\frac{k e^2}{2},
\]
which gives
\begin{equation}
    \Eprivate=\frac{m\beta(1-h)}{k}.
    \label{eq:ep}
\end{equation}
The coordinated system chooses effort to maximize operating surplus across both routes net of effort cost, yielding
\begin{equation}
    \Ecoordinated
    =\frac{\beta\left[(p-c_S)(1-h)+(p-c_R)h\right]}{k}.
    \label{eq:ec}
\end{equation}

Social welfare includes consumer surplus and real production/distribution surplus. Since demand is $d=a+\beta e-p$, aggregate consumer surplus is $d^2/2$. Maximizing total surplus gives
\begin{equation}
    \Esocial
    =\frac{\beta\left[a-c_S(1-h)-c_Rh\right]}{k-\beta^2}.
    \label{eq:ew}
\end{equation}
The restriction $k>\beta^2$ guarantees strict concavity of the social problem.

\begin{theorem}[Effort ordering]\label{thm:welfareordering}
Under $a>p>c_S>c_R\geq0$, $0<m<p-c_S$, $\beta>0$, $k>\beta^2$, and $0<h<1$,
\begin{equation}
    \Eprivate<\Ecoordinated<\Esocial.
    \label{eq:ordering}
\end{equation}
\end{theorem}

\begin{proof}
Using \eqref{eq:ep}--\eqref{eq:ec},
\[
\Ecoordinated-\Eprivate
=\frac{\beta}{k}\left[(p-c_S-m)(1-h)+(p-c_R)h\right]>0,
\]
because $p-c_S-m>0$ and $p-c_R>0$. Let
\[
M=(p-c_S)(1-h)+(p-c_R)h>0.
\]
The numerator in \eqref{eq:ew} equals $M+(a-p)$. Hence
\[
\Esocial-\Ecoordinated
=\frac{\beta\left[k(a-p)+\beta^2M\right]}{k(k-\beta^2)}>0,
\]
because $a>p$, $M>0$, and $k>\beta^2$. This proves the strict ordering.
\end{proof}

The first wedge reflects incomplete appropriability within the vertical system: $S$ captures only its own margin rather than the system margin associated with product improvement. The second wedge reflects the additional consumer-surplus value of quality improvement. These are distinct objects; vertical coordination does not generally implement the social optimum.

\subsection{Widening wedges under reallocation}

Differentiating the private--coordinated gap gives
\begin{equation}
    \frac{d(\Ecoordinated-\Eprivate)}{d\eta}
    =\frac{\beta h'(\eta)}{k}\left[(c_S-c_R)+m\right]>0.
    \label{eq:wedgepc}
\end{equation}
The gap widens because reallocation simultaneously reduces the information producer's private footprint and shifts system activity toward the lower-real-cost route.

The private--social gap also widens. Differentiating \eqref{eq:ew} and \eqref{eq:ep} gives
\begin{equation}
    \frac{d(\Esocial-\Eprivate)}{d\eta}
    =\beta h'(\eta)
    \left[
    \frac{c_S-c_R}{k-\beta^2}+\frac{m}{k}
    \right]>0.
    \label{eq:wedgepw}
\end{equation}
Thus downstream reallocation can increase the social value of information provision at precisely the time when the decentralized information producer has a weaker incentive to provide it.

\subsection{Organizational implication}

The benchmark also supports a limited organizational implication. Consider a separate retention decision in which maintaining the information-producing downstream capacity is privately worthwhile only above a threshold. At an interior private retention-indifference point, the manufacturer's incremental private return from retention is zero. If the retained information-producing node earns a positive equilibrium payoff and the associated product improvement raises consumer surplus, social welfare strictly favors retention at that point. By continuity, provided there is a genuine private non-retention side of the threshold, there exists a nonempty local region in which private organization drops the information-producing capacity while social welfare favors retaining it.

This is deliberately a local and conditional result. It is not a general foreclosure theorem, nor does it imply that retaining the information-producing route is always efficient. Its role is to translate the effort wedge into an organizational margin without adding a new strategic mechanism.
